\section{Versuche}

\subsection{Versuch 1: pH-Wert Abschätzung und Bestimmung}
\subsubsection{Aufgabenstellung}
Es sollen  pH-Werte angegebener Lösungen geschätzt und bestimmt werden, um festzustellen ob diese in Wasser sauer, neutral oder alkalisch reagieren.
\subsubsection{Durchführung}
Neun Reagenzgläser wurden mit \SI{0.1}{\frac{\mol}{l}} Salzsäure, \SI{0.01}{\frac{\mol}{l}}, Salzsäure, demineralisiertem Wasser, Natriumcarbonatlösung, Natriumhydrogencarbonatlösung, Natriumacetatlösung, Kaliumnitratlösung, Ammoniumchloridlösung, Hydroxylammoniumchloridlösung bis zur Hälfte befüllt. Die genannten Lösungen wurden aus den bereitgestellten Chemikalien hergestellt. Zuerst wurden die pH-Werte geschätzt, anschließend mit Indikatorpapier bestimmt und abschließend mit einem pH-Meter gemessen.

\subsubsection{Beobachtungen}
\begin{table}[H]
    \centering
    \caption{}
    \label{tab:theoretische-werte}
    \begin{tabular}{ c|cc  }
        \hline
        $H_{x}$-Linie & $n_2$ & $\lambda_{\text{theo}}$ in \si{nm} \\
        \hline
        $H_{\alpha}$  & 3     & 656                                \\
        $H_{\beta}$   & 4     & 486                                \\
        $H_{\gamma}$  & 5     & 433                                \\
        \hline
    \end{tabular}
\end{table}

\subsubsection{Auswertung}
\subsubsection{Fehlerbetrachtung}


\subsection{Versuch 2: Verwendung geeigneter Indikatoren}
\subsubsection{Aufgabenstellung}
\subsubsection{Durchführung}
\subsubsection{Beobachtungen}
\subsubsection{Auswertung}
\subsubsection{Fehlerbetrachtung}


\subsection{Versuch 3: Titration von Essigsäure}
\subsubsection{Aufgabenstellung}
\subsubsection{Durchführung}
\subsubsection{Beobachtungen}
\subsubsection{Auswertung}
\subsubsection{Fehlerbetrachtung}

\subsection{Berechnung einer Puffermischung}
\subsubsection{Aufgabenstellung}
\subsubsection{Berechnung}

\newpage
\section{Zusammenfassung}
